This thesis introduced an extension to the Vessim framework \cite{vessim}, an already existing co-simulation testbed for carbon-aware applications and systems, to allow for better simulations of battery packs, depending on the experimentation context. The extension included the definition of an energy storage interface, allowing the implementation of models for battery systems of different complexity, and established policies to be used for energy management during the simulation. Additionally, the improved architecture allowed \textit{visibility} and \textit{control} over the energy storage system through controllers and message passing, using the underlying scheduler of Mosaik \cite{mosaik}.

To evaluate the usefulness of utilizing various models for simulating lithium-ion battery packs, this work conducted an analysis, integrating a total of four models into the described framework. The first model described a simple, lossless battery that did not consider any power limitations, whereas the second implementation was based on the C-L-C model by Kazhamiaka et al. \cite{kazhamiaka2019tractable}, adding inefficiencies and power constraints. The third model made use of the PyBaMM framework \cite{pybammbatteries} to simulate an individual cell using the physics-based Single Particle Model, which was scaled to represent a battery pack, before the final model applied multiple of these electrochemical models together with an external circuit to reproduce the behavior of a realistic battery pack.

The experiments showed that, despite its abilities to determine overall trends for a battery's state and exchanged grid energy, the lack of accuracy by the simplest model severely limited its usability in experiments studying energy management of computing infrastructure. Adding linear constraints and inefficiencies with the C-L-C model massively improved such estimations with only minimal performance overhead, making those kinds of models ideal for most scenarios, due to its simple process of parameterization.

Although SPMs offered a more accurate insight into the electrochemical processes of battery cells, the high computational cost for stepping the models and the complexities to determine an accurate parameterization made them impractical for short-term scenarios. They could, however, still be used to study long-term effects like battery degradation, although further research would have to be conducted. Modeling a whole battery pack instead of a single scaled-up cell was found to be inefficient, as only minor power-loss effects are noticeable for small battery packs, and larger battery packs come with significant execution times, raising questions about the utility of such detailed modeling.

To extend on this thesis, future work should focus on benchmarking the used simulation tools against real systems. For instance, there is still much uncertainty, whether co-simulation testbeds such as Vessim accurately determine the power-deltas that are used to (dis-)charge the energy storage components compared to real-world power systems. On top, there is a need to evaluate models like the SPMs by PyBaMM, and a parameterization of an individual cell against real-world data to ensure accurate long-term behavior. Kazhamiaka et al. \cite{kazhamiaka2019tractable} also described other tractable models for lithium-ion battery cells with more advanced approximations like the C-L-L or the L-L-Q model to be investigated in the same ways as described in this thesis.